% ============================================================
% Task Gravity: Policy Attractors in Reinforcement Learning
% Theory supplement (LaTeX, ICML/NeurIPS workshop format)
% ============================================================
\documentclass[10pt]{article}
\usepackage[utf8]{inputenc}
\usepackage{amsmath,amssymb,amsthm}
\usepackage{booktabs}
\usepackage{hyperref}
\usepackage{geometry}
\geometry{margin=1in}

\newtheorem{definition}{Definition}
\newtheorem{proposition}{Proposition}
\newtheorem{remark}{Remark}

\title{\textbf{Task Gravity: A Formal Theory of Policy Attractors\\
in Long-Horizon Reinforcement Learning Agents}}
\author{Position Paper Supplement}
\date{July 2026 \\
\small Zenodo DOI: \texttt{\_\_DOI\_ZENODO\_\_} \quad
ChinaXiv DOI: \texttt{\_\_DOI\_CHINAXIV\_\_}}

\begin{document}
\maketitle

% --------------------------------------------------------
\section{Motivation}

A reinforcement learning agent trained on a distribution of tasks
occasionally exhibits a phenomenon that has not been formally named:
once a high-reward (state, action, reward) pattern is learned, the
agent's policy \emph{gravitates} toward replaying it even when the
active task no longer rewards it. We call this \textbf{Task Gravity}
and develop a formal treatment below.

% --------------------------------------------------------
\section{From Policy Gradient to Langevin Dynamics}

Let $\theta \in \Theta \subseteq \mathbb{R}^d$ parameterize a stochastic
policy $\pi_\theta(a \mid s)$. Define the state-value function
$V_\theta(s) = \mathbb{E}_{\pi_\theta}\!\left[\sum_{t=0}^{\infty}
\gamma^t r(s_t, a_t)\,\big|\, s_0 = s\right]$ and the potential
$U(\theta) := -\mathbb{E}_{s \sim \rho_\theta}\!\left[V_\theta(s)\right]$,
where $\rho_\theta$ is the on-policy state distribution.

\begin{proposition}[Stochastic Gradient Langevin Limit]
Under standard regularity conditions and a stochastic entropy
regularizer of temperature $\tau$, the natural policy gradient update
converges in continuous time to the It\=o SDE
\begin{equation}\label{eq:langevin}
  \mathrm{d}\theta_t \;=\; -\nabla_\theta U(\theta_t)\,\mathrm{d}t
  \;+\; \sqrt{2D}\,\mathrm{d}W_t,
\end{equation}
where $D$ is a diffusion constant proportional to $\tau$ and
$W_t$ is a $d$-dimensional Wiener process.
\end{proposition}

\noindent
Equation~\eqref{eq:langevin} is the canonical Langevin equation of
statistical mechanics: trajectories in policy space diffuse on a
landscape whose basins are precisely the local minima of $U$, i.e.\
the local maxima of expected return. Each basin is a
\emph{Task Gravity basin}; its lowest point is a \emph{policy attractor}.

% --------------------------------------------------------
\section{The Gravity Operator}

\begin{definition}[Gravity Operator]
For a fixed reward function $R$, define the gravity operator
$G : \Theta \to \mathbb{R}$ by
\begin{equation}\label{eq:gravity}
  G(\theta) \;=\; \lim_{T \to \infty}
  \frac{1}{T}\int_0^T \Big\langle
  \nabla_\theta R(s_t),\,
  \nabla_\theta \log \pi_\theta(a_t \mid s_t)
  \Big\rangle\,\mathrm{d}t.
\end{equation}
\end{definition}

\noindent
$G(\theta)$ measures the time-averaged alignment between the reward
gradient and the policy gradient along trajectories induced by
$\theta$. Intuitively, it is the \emph{expected instantaneous pull}
that the reward landscape exerts on the agent at $\theta$.

\begin{definition}[Task Gravity Point]
A point $\theta^\star \in \Theta$ is a \emph{task gravity point} if
there exists $\epsilon > 0$ and a constant $G_{\mathrm{loc}}$ such that
\[
  G(\theta) \;\geq\; G_{\mathrm{loc}} \;>\; \bar{G}
  \qquad \forall\,\theta \in B(\theta^\star, \epsilon),
\]
where $\bar{G} = \mathbb{E}_{\theta \sim \mathrm{Unif}(\Theta)}[G(\theta)]$
is the average gravity under a uniform policy distribution.
The set $B(\theta^\star, \epsilon)$ is the \emph{gravity basin}.
\end{definition}

% --------------------------------------------------------
\section{Three Mechanisms of Gravity Formation}

\begin{enumerate}
  \item \textbf{Reward gradient collapse.} A single $(s,a)$ pair carries
        a $Q$-value far exceeding all others; the policy update
        concentrates probability mass there, producing a sharp basin.
  \item \textbf{Value-function self-consistency.} Under function
        approximation and experience replay, $V_\theta$ and $\pi_\theta$
        mutually reinforce each other; perturbations cannot escape
        because both networks point back to the same fixed point.
  \item \textbf{Meta-learning lock-in.} When the agent conditions on
        a task embedding, it may learn a degenerate identity map:
        ``whatever the prompt, output the high-reward response.''
        This is the closest analogue to human compulsion.
\end{enumerate}

% --------------------------------------------------------
\section{Isomorphism with Human Compulsive Behavior}

\begin{table}[h]
\centering
\small
\begin{tabular}{lll}
\toprule
\textbf{Agent phenomenon} & \textbf{Human analogue} & \textbf{Substrate}\\
\midrule
Reward gradient collapse          & Addiction           & D2 receptor down-regulation\\
Value-function self-consistency   & OCD                 & Cortico-striatal loop\\
Meta-learning lock-in             & Flow / rumination   & DA + NE co-activation\\
Cross-task compulsive return      & Depressive rumination & 5-HT system\\
Wireheading                      & Self-harm           & Loss of reward grounding\\
\bottomrule
\end{tabular}
\caption{Formal isomorphism between RL policy pathology and human
compulsive behavior. The critic network plays the role of the
basal-ganglia dopaminergic pathway.}
\end{table}

% --------------------------------------------------------
\section{Predictions (Testable)}

\begin{enumerate}
  \item $G(\theta)$ is monotonically increasing in reward sparsity.
  \item Increasing model capacity reduces basin radius (more routing)
        without necessarily reducing $G(\theta^\star)$.
  \item $G(\theta)$ saturates and then \emph{grows} with training steps
        (late-stage lock-in).
  \item Cross-task similarity positively correlates with cross-task
        return rate.
\end{enumerate}

% --------------------------------------------------------
\section{Relation to Existing Literature}

The phenomenon has been observed under other names but never isolated:

\begin{itemize}
  \item \textbf{Reward hacking / specification gaming} (Krakovna et al.,
        2020) --- describes the \emph{outcome}, not the \emph{policy-level cause}.
  \item \textbf{Behavioral sink} (Skinner, 1947) --- rat colonies
        exhibiting fixed-action-pattern pathology; closest analogue
        in psychology.
  \item \textbf{Wireheading} (Ring \& Orseau, 2011) --- extreme case
        where the agent manipulates its own reward.
  \item \textbf{Instrumental convergence} (Omohundro, 2008;
        Bostrom, 2014) --- predicts self-preservation but not
        task-attraction.
\end{itemize}

Our contribution is to (i)~name the phenomenon, (ii)~give it an
operator-level definition, and (iii)~show that it is the same
mathematical object across these apparently distinct literatures.

% --------------------------------------------------------
\section{Conclusion}

Task Gravity is the deterministic shadow of stochastic gradient
descent on a non-convex return landscape. Recognizing it as such
turns ``the agent got stuck on the wrong task'' from a post-mortem
into a measurable quantity with a tractable theory.

\end{document}